\begin{table}[htbp]
\centering
\begin{tabular}{|l|c|c|c|c|c|}
    \hline
    Signal & $n_{\mathrm{SVJ}}^{\mathrm{PN}} = 0$ & $n_{\mathrm{SVJ}}^{\mathrm{PN}} = 1$ & $n_{\mathrm{SVJ}}^{\mathrm{PN}} = 2$ & $n_{\mathrm{SVJ}}^{\mathrm{PN}} \geq 3$ & Inclusive \\
    \hline
    $m_{\Phi}=600$ GeV, $r_{\mathrm{inv}}=0.3$ & 0.30\% & 0.28\% & 0.09\% & 0.02\% & 0.68\% \\
    $m_{\Phi}=800$ GeV, $r_{\mathrm{inv}}=0.3$ & 0.66\% & 0.65\% & 0.22\% & 0.03\% & 1.54\% \\
    $m_{\Phi}=1000$ GeV, $r_{\mathrm{inv}}=0.3$ & 0.64\% & 0.84\% & 0.33\% & 0.03\% & 1.84\% \\
    $m_{\Phi}=1500$ GeV, $r_{\mathrm{inv}}=0.3$ & 0.64\% & 0.95\% & 0.36\% & 0.04\% & 2.00\% \\
    $m_{\Phi}=2000$ GeV, $r_{\mathrm{inv}}=0.1$ & 0.08\% & 0.23\% & 0.16\% & 0.03\% & 0.50\% \\
    $m_{\Phi}=2000$ GeV, $r_{\mathrm{inv}}=0.3$ & 0.37\% & 0.64\% & 0.30\% & 0.04\% & 1.35\% \\
    $m_{\Phi}=2000$ GeV, $r_{\mathrm{inv}}=0.5$ & 0.57\% & 0.64\% & 0.21\% & 0.03\% & 1.44\% \\
    $m_{\Phi}=2000$ GeV, $r_{\mathrm{inv}}=0.7$ & N/A & N/A & N/A & N/A & N/A \\
    $m_{\Phi}=3000$ GeV, $r_{\mathrm{inv}}=0.3$ & 0.19\% & 0.42\% & 0.27\% & 0.05\% & 0.93\% \\
    $m_{\Phi}=4000$ GeV, $r_{\mathrm{inv}}=0.3$ & N/A & N/A & N/A & N/A & N/A \\
    \hline
\end{tabular}
\caption{Fraction of \emph{generated} Run 2 signal events landing in the ABCD signal region A ($N_A / N_{\mathrm{generated}}$), inclusively and in bins of $n_{\mathrm{SVJ}}^{\mathrm{PN}}$. $N_{\mathrm{generated}}$ is the Run2-combined, luminosity-scaled yield before any selection, read from each signal sample's CutFlow/Initial entry -- the same normalization NMinusOne\_maker\_RA2.py uses for its Lund-plane correction. Unlike the region-A-only table above, these nSVJ columns share a common denominator and sum to the Inclusive column.}
\label{table:signal_region_a_vs_generated}
\end{table}
